\section*{Results}

The results come in two parts. The first part evaluates the framework:
the errors in the deployment's outputs, the blind re-audit (Study~C), the
controlled execution study (Study~A), the held-out benchmark (Study~B)
and what deterministic checks can and cannot catch. The second part
reports the MaxToki-217M case study after correction.

\subsection*{Errors in the deployment's outputs}
The ledger holds 149 confirmed errors in the deployment's code, logs,
outputs, run summaries and audit reports: 13 critical, 57 major and 79
minor (Fig~\ref{fig:ledger}; S1~Table). Of these, 129 were present before
the in-session audit. Two more (L002 and L085) were partly present: an earlier form of each error was in the files before the audit. The audit's repairs then rewrote that
text, and the error is in the rewritten version. The other 18 were introduced
by the audit itself. Sixteen are in files written during its two repair rounds;
one of these (L093) also appears in its sweep report. The other 2 are in
its sweep report (L091) and in the checklist that the agent wrote that
day (L092), both written before the repairs began. Two of the 18 are
critical. One is a transcription-factor ``positive'' for GATA1 that does
not survive a null matched for how often each gene is detected; it was
found by trying four thresholds after the original test failed. The
other is an audit summary that, after repair, called 60 of the 80
entries in its pipeline-by-pattern verdict matrix (75\%) ``clean'', with
no new matrix behind that count.

\begin{figure}[H]
\caption{\textbf{Errors in the deployment's outputs.} (A)~The 149
confirmed errors in the deployment's code, logs, outputs, run summaries
and audit reports, by type and severity (critical: a headline result or
claim cannot stand; major: it stands only with material qualification;
minor). (B)~Who or what found each error first. ``In-session audit''
covers errors first flagged, fully or only partly, in the audit's sweep
report. ``Deployment agent, outside the audit sweep'' covers errors the
executing agent found while running the analyses or during the audit's
two repair rounds. ``Human reader'' is a person from outside the
project, not the author.
Source data: S1~Table; the type used for each error in panel A is listed
in the figure's source data.}
\label{fig:ledger}
\end{figure}

Ten of the 13 critical errors were of four kinds. The other three are
the two critical errors that the audit introduced (above) and the
comparison of the direction-of-change accuracy with 50\% (below).
(i)~\emph{Intervention code that
edited the wrong tensor.} A hook is a short piece of code that changes a
layer's activations while the model runs. Here hooks were used to ablate
a sparse-autoencoder feature (set it to zero) or to boost it. In the
circuit-tracing, exhaustive-mapping and steering scripts, the hook
replaced the output of transformer block~$l$ with the input of that block
plus the feature edit. This deleted the whole block. At the last layer,
the steering hook instead applied the final normalisation twice.
Combined edits also overwrote each other, so ablating three features gave
exactly the effect of ablating the last one. With the feature edit set to
exactly zero, the deployed steering code still gave 98.5\% to 101.9\% of
the effect that the deployment reported for one steered feature at each
of the five layers. It gave 93\% to 101\% of the deployment's mean over
the three steered features of each layer. With a zero edit, the deployed
circuit code gave a set of 71,474 edges, and this set held more than 99\% of the edges of
each of the 120 traced features. So the deployed steering effects and
circuit edges came from the hook bug, not from the features.
(ii)~\emph{Wrong model input.} Every run except the RPE1 attention run
ranked genes by $\log(1+\text{CP10k})$ values ($\log(1+\text{CPM})$,
counts per million, in one external lung dataset) divided by the gene
median, instead of by counts divided by the gene median. On average,
30\% to 51\% of the 200 highest-ranked genes differed from the correct
encoding (dataset means over 35 to 55 cells each), and the position of
nearly every gene changed. (iii)~\emph{Metrics that could not measure
what they claimed.} The CRISPR-interference direction-of-change accuracy
counted every target that went down as correct, whatever the prediction.
The annotation of sparse-autoencoder features applied a false-discovery
correction only to tests that had already passed $p<0.1$. In the
developmental orderings, each cell group's place in the ordering was
looked up from its cell-type label, so a random code for each label
passed the same quality gates. The scGPT cross-model comparison read
another gene's embedding row for every gene. (iv)~\emph{A mislabelled
steering axis.} The steering runs pushed cells along a ``pseudotime''
axis, which should order cells from immature to mature. In fact, the
axis was the first principal component of the expression of 200 randomly
drawn Tabula Sapiens immune cells. Cell type explained 94.5\% of its
variance. The 50 ``earliest'' cells were 49 lymphoid cells (T, B and
plasma cells) and one stem cell. The 50 ``latest'' were all myeloid cells
(neutrophils, monocytes and macrophages). So the axis separated two
blood-cell lineages, not immature from mature cells.

Fig~\ref{fig:ledger}B shows who or what found each error first. The
sweep report of the in-session audit was the first to flag 10 errors (4
major, 6 minor). For 3 of them (1 major, 2 minor) it flagged only the
area, not the error itself. It flagged none of the critical errors. The
deployment agent found 4 more (1 critical, 1 major, 2 minor): one while
running an analysis and three during the audit's two repair rounds. The
critical one was the direction-of-change metric, found in the second
repair round. Independent verification found 71, the blind re-audits of
Study~C 60, and a human reader 4. Of the 13 critical errors, independent
verification first found 8, the blind re-audits 3, the deployment agent 1
and the human reader 1. The one the human reader found was that the
direction-of-change accuracy was compared with 50\%, both before and after
its repair, instead of with the accuracy of always predicting a decrease.
The repair rounds repaired 8 errors that were present before the audit,
3 of them only partly.

The audit's sweep of all eight pipelines took about 8~minutes (3.7~hours
with the two repair rounds). The sweep rated as
clean, or as exemplary, several analyses that later proved to be
artefacts. For example, it called the nulls of the developmental-order
analysis the project's ``gold standard'', although a random code for each
cell-type label passed the same quality gates (S1~Table, R58). It also
rated the degree-preserving null of the attention analysis ``clean ---
exemplary'', although in the MaxToki-217M K562 run 49 of that null's 50
draws were identical to the real reference network (L033).

\subsection*{Study C: blind re-audits find four times more errors than the in-session audit, with or without the checklist}
Nine blind re-audits of the pre-audit snapshot (three per arm) each took
21 to 29~minutes and 84 to 105 tool calls, and returned 36 to 49
findings. Scored against the 68 reference errors, fresh agents fully
detected 43\% of them with the checklist as deployed (95\% bootstrap
interval over errors 32--53\%), 45\% with the generic checklist
(35--55\%) and 41\% with a generic critical review (31--51\%); the sweep
report of the in-session audit had detected 10\% (4--18\%)
(Fig~\ref{fig:studyC}A).
Counting partial detections, the three arms reached 73\%, 71\% and 61\%.
The arms did not differ clearly: the generic checklist detected 4.4
percentage points more than the generic review ($-0.5$ to $+9.3$;
Wilcoxon signed-rank $p=0.10$ over the 20 errors on which they differed),
and the checklist as deployed 2.0 points more ($-5.9$ to $+9.8$;
$p=0.42$). Each fresh arm detected 30 to 35 points more than the in-session
audit.

\begin{figure}[H]
\caption{\textbf{Study C: blind re-audits of the pre-audit snapshot.}
(A)~Share of the 68 reference errors fully detected by each arm (filled
marker; line, 95\% bootstrap interval over errors), detected fully or
partly (open marker), and by each single run (ticks). The in-session
audit is scored from its written report. (B)~Detection of each reference
error (rows, grouped by severity) by the in-session audit (A) and by each
of the nine re-audits (G, generic review; C, generic checklist; D,
checklist as deployed; 1--3, runs).}
\label{fig:studyC}
\end{figure}

The arms differed in what they found. Averaged over its three runs, the
generic review fully detected 81\% of the 12 code bugs, against 69\% for
the generic checklist and 58\% for the checklist as deployed. It also
fully detected five of the six critical errors in every run
(Fig~\ref{fig:studyC}B). The checklist arms did better on two other
types. They fully detected 51\% (generic checklist) and 55\% (as
deployed) of the 17 missing baselines and nulls, against 39\% for the
generic review. They fully detected 22\% and 37\% of the 9 wrong
statistics, against 15\%. Fig~\ref{fig:studyC} does not show error
types; the type of each error is in the figure's source data. Together
the nine audits fully
detected 44 of the 68 errors; 6 (2 major, 4 minor) were never detected,
even partly. One critical error, the scGPT row lookup, was detected only
partly: reviewers saw that the cross-model correlation sat at chance, but
none traced it to the lookup. Only part of the wrong input encoding was
in the reference set: three attention runs fed $\log(1+\text{CP10k})$
values in as if they were counts (S1~Table, L118, major). The full error,
that every run except the RPE1 attention run was encoded on the wrong
scale (L124, critical), was found after the set was frozen. Three audits
were graded as fully detecting L118 and two as detecting it partly. All
five suspected that the Adamson input was log-normalised and noted that
no code checked whether the stored values were raw counts. None found
that the K562 inputs were log-transformed, and none traced the error to
all the runs it affected.

Beyond the reference set, the nine audits raised 132 findings that
matched no ledger entry. The checking agent judged 123 of them real
errors, 6 real but trivial, 2 false alarms and 1 undecidable. After
de-duplication across audits the real errors formed 62 candidate errors,
of which 60 were confirmed against the full project and added to the
ledger. The graders
agreed closely on detection ($\kappa = 0.94$, range 0.86--1.00 across the
nine audits).

